\subsection{Convexity, inflection, and acceleration}

The second derivative measures how the first derivative changes.

\begin{theorembox}[Shape from the second derivative]
If $f''(x)>0$ on an interval, the graph is convex there. If $f''(x)<0$, the graph is concave there. A change of sign of $f''$ indicates a possible inflection point.
\end{theorembox}

\begin{examplebox}[Shape of a cubic]
For
\[
f(x)=x^3-3x,
\]
we have
\[
f''(x)=6x.
\]
Thus the graph is concave for $x<0$, convex for $x>0$, and changes convexity at the origin.
\end{examplebox}

\begin{examplebox}[Acceleration]
If position is
\[
s(t)=t^3-3t,
\]
then velocity and acceleration are
\[
v(t)=s'(t)=3t^2-3,
\qquad
a(t)=s''(t)=6t.
\]
The sign of $a(t)$ tells us whether velocity is increasing or decreasing.
\end{examplebox}

\begin{interpretationbox}[One calculation, two languages]
For a graph, the second derivative describes changing slope and convexity. For motion, it describes changing velocity and acceleration. The mathematics is the same; the interpretation depends on the model.
\end{interpretationbox}
